%====================
% SELECTED PROJECTS
% KONUX: lead with the fine-tuning project -- it is the one that shows models
% trained, benchmarked against baselines and selected on measured results, with
% data-preparation / training / evaluation pipelines and experiment tracking.
% The RAG QA bot follows.
%====================
\documentclass[letter,10pt]{article}
\usepackage{TLCresume}
\usepackage{hyperref}
\begin{document}

\subsection{Vision-Language Model Fine-tuning – Image Captioning (Freelance) \hfill \href{https://huggingface.co/sunil448832/blip-model}{HuggingFace Model}}
\begin{zitemize}
\item Fine-tuned \textbf{LLaVA 1.6 (Mistral-7B)} with \textbf{LoRA} (r=128, $\alpha$=256) via \textbf{DeepSpeed ZeRO-3} for domain-specific image captioning, then \textbf{benchmarked} it against \textbf{CLIP} similarity and a fine-tuned \textbf{BLIP} baseline to select the best-performing approach on measured results.
\item Built end-to-end \textbf{data preparation, training, and evaluation} pipelines with \textbf{Weights \& Biases} experiment tracking; published the fine-tuned \textbf{BLIP} model to the \textbf{HuggingFace Hub}.
\item \textbf{Stack:} PyTorch, HuggingFace Transformers, LLaVA, DeepSpeed, LoRA, CLIP, BLIP, W\&B.
\end{zitemize}

\subsection{QA Bot for PAN Card Queries (RAG) \hfill \href{https://github.com/sunil448832/Natural-Language-Processing/tree/master/RAG(Retrieval\%20Augment\%20Generation)}{GitHub Repository}}
\begin{zitemize}
\item Developed a \textbf{RAG} pipeline for PAN card query answering over a custom document corpus, using \textbf{Sentence Transformers} (MiniLM-L6-v2) embeddings with \textbf{FAISS} vector retrieval.
\item Queried a \textbf{quantized Mistral-7B} for contextual response generation; deployed via \textbf{FastAPI}.
\item \textbf{Stack:} PyTorch, FAISS, HuggingFace, FastAPI.
\end{zitemize}

\end{document}
